% Komunikačné protokoly v IoT
\section{Komunikačné protokoly}

V tejto kapitole sú podrobne rozpracované protokoly a komunikačné mechanizmy, ktoré sú bežne použité v IoT systémoch. Výber konkrétneho protokolu závisí od požiadaviek aplikácie (energetická náročnosť, latencia, spoľahlivosť, šírka pásma, bezpečnosť a škálovateľnosť). Nižšie sú popísané vlastnosti, prevádzkové režimy a návrhové kompromisy pre vybrané protokoly.

\subsection{MQTT}

MQTT (Message Queuing Telemetry Transport) je ľahký publish/subscribe protokol postavený nad TCP/IP, optimalizovaný pre nízku šírku pásma a obmedzené zariadenia. Architektúra zahŕňa vydavateľov (publishers), odoberateľov (subscribers) a sprostredkovateľa (broker).

Kľúčové vlastnosti:
\begin{itemize}
  \item QoS (Quality of Service) úrovne: 0 (najneskôr raz), 1 (aspoň raz), 2 (presne raz) — každá úroveň prináša rozdielnu réžiu a zložitosť implementácie.
  \item Retained správy: broker môže uchovávať poslednú hodnotu pre dané téma a odoslať ju novým odberateľom.
  \item Last Will and Testament: umožňuje signalizovať náhly odchod klienta.
  \item Prenos nad TCP s voliteľným TLS pre zabezpečenie šifrovania a overenia identity.
\end{itemize}

Návrhové úvahy: MQTT je vhodný pre telemetriu a scenáre so sporadickými aktualizáciami; nedostatky zahŕňajú závislosť na spojenom TCP kanáli a možnú potrebu správy brokerov pre škálovanie.

\subsection{CoAP}

CoAP (Constrained Application Protocol) je protokol navrhnutý pre obmedzené zariadenia a siete, pracujúci nad UDP a sledujúci štýl REST (representational state transfer). Jeho základné mechanizmy zahŕňajú správy typu Confirmable (potvrdzované) a Non-confirmable, a podporu pozorovania (Observe) proaktívne zasielajúce aktualizácie.

Kľúčové vlastnosti:
\begin{itemize}
  \item Nízka režijná réžia vďaka UDP a kompaktnému formátu hlavičky.
  \item Podpora metód podobných HTTP: GET, POST, PUT, DELETE.
  \item Možnosť zabezpečenia pomocou DTLS (Datagram TLS) alebo OSCORE (Object Security for Constrained RESTful Environments).
\end{itemize}

Návrhové úvahy: CoAP je vhodný tam, kde je potrebná REST‑ová interakcia pri nízkej režii, napríklad pri komunikácii medzi senzorom a lokálnym bránovým zariadením. Vzhľadom na použitie UDP je potrebné riešiť spoľahlivosť a poradie správ aplikačne tam, kde to záleží.

\subsection{HTTP/REST a WebSocket}

HTTP/REST je štandardný protokol pre webové služby a často sa používa v IoT pri integrácii so servermi a cloudovými rozhraniami. Nevýhodou je vyššia režijná záťaž (TCP, hlavičky) v porovnaní s MQTT a CoAP.

WebSocket poskytuje perzistentné obojsmerné spojenie nad TCP, vhodné pre real‑time dashboardy a ovládanie zariadení s nízkou latenciou. Pre nízkoenergetické zariadenia sú tieto protokoly menej vhodné bez proxy alebo brány, ktorá agreguje komunikáciu.

\subsection{Narrowband a LPWAN technológie: LoRaWAN, NB‑IoT, Sigfox}

LoRaWAN je sieťová vrstva postavená nad moduláciou LoRa, vhodná pre dlhý dosah a nízku spotrebu energie. Podporuje tri triedy zariadení (A,B,C) s rôznymi režimami prijímania, adaptívne riadenie výkonu (ADR) a režimy prispôsobenia prenosovej rýchlosti.

NB‑IoT (Narrowband IoT) a Sigfox sú sieťové technológie poskytované mobilnými/operátorskými infraštruktúrami s dôrazom na nízku spotrebu, dlhú životnosť batérií a veľké pokrytie. NB‑IoT poskytuje lepšiu priepustnosť a integráciu s mobilnými sieťami než Sigfox.

Návrhové úvahy: LPWAN technológie sú vhodné pre scenáre s nízkou frekvenciou zasielania údajov a veľkým dosahom, nie sú vhodné pre aplikácie s vysokými nárokmi na latenciu alebo vysokú priepustnosť.

\subsection{Bluetooth Low Energy (BLE) a Zigbee}

BLE poskytuje nízku spotrebu pri krátkom dosahu a model GATT (Generic Attribute Profile) na organizáciu služieb a charakteristík. BLE je vhodné pre nositeľné zariadenia a lokálne senzory, pričom nové verzie zlepšujú prenosovú rýchlosť a energetickú efektívnosť.

Zigbee je mesh‑orientovaná sieťová technológia bežne používaná v domácich automatizáciách; podporuje samoliečenie siete a distribuované smerovanie. Oba štandardy majú vlastné bezpečnostné mechanizmy (pairing, link keys), ktoré je potrebné správne implementovať.

\subsection{6LoWPAN, Thread a IPv6}

6LoWPAN umožňuje prenos IPv6 paketov cez nízkopásmové rádiové siete (IEEE 802.15.4), čím poskytuje možnosť priamej adresovateľnosti zariadení a používaniu štandardných protokolov IPv6. Thread stavia na 6LoWPAN a poskytuje mesh‑funkcionalitu s bezpečnosťou na sieti.

\subsection{Bezpečnostné aspekty protokolov}

Pri navrhovaní IoT riešení je rozhodujúce zabezpečiť ochranu dát a integritu komunikácie. Základné odporúčania:
\begin{itemize}
  \item Používať šifrovanie vrstvy transportu (TLS/DTLS) alebo aplikačné zabezpečenie (OSCORE).
  \item Overovať a autorizovať zariadenia pomocou silných kľúčov alebo certifikátov; vyhýbať sa statickým, nezabezpečeným kľúčom v zariadení.
  \item Implementovať aktualizácie firmvéru zabezpečeným spôsobom (signed firmware, over‑the‑air aktualizácie s overením integrity).
\end{itemize}

\subsection{Kritériá výberu protokolu}

Pri rozhodovaní o voľbe protokolu je vhodné zvážiť nasledujúce kritériá: energetická náročnosť, spoľahlivosť (QoS), latencia, dostupná šírka pásma, režijné náklady, požadovaná úroveň zabezpečenia, topológia siete a požiadavky na škálovanie.

V nasledujúcich kapitolách sú tieto protokoly demonštrované na konkrétnych príkladoch a implementáciách na platforme ESP32 a prostredí MicroPython.
